\chapter{Conclusioni e Sviluppi Futuri}
Questo lavoro di tesi si è posto l'obiettivo di analizzare l'affidabilità dei Large Language Models nel dominio della Vulnerability Detection, approfondendo non solo la loro suscettibilità agli attacchi testuali ma mappando il processo di tale vulnerabilità attraverso la metodologia della Mechanistic Interpretability.

L'analisi condotta sulla baseline ha evidenziato come le architetture odierne, seppur addestrate specificamente per la programmazione come le varianti coder, non possiedono ancora le capacità necessarie per operare in contesti di sicurezza senza la supervisione umana. Questa fragilità è stata confermata dai risultati degli attacchi Black-Box: modelli adottati in ambienti lavorativi presentano un alto tasso di suscettibilità alle perturbazioni introdotte.
Le metodologie di attacco presentate nello studio hanno portato alla luce diverse mancanze dei modelli, possibilmente fatali, registrando un tasso di successo molto alto.
Attacchi come il Prompt Injection basato sull'inserimento di perturbazioni testuali, anche sintatticamente slegate dalla logica del codice sorgente analizzato, sono sufficienti a sovrascrivere le direttive di sistema e a forzare una classificazione errata.

Il contributo principe di questa ricerca risiede tuttavia nell'esplorazione vettoriale del Residual Stream. Tramite il processo di estrazione implementato, è stato possibile superare la pura analisi comportamentale di input-output e visualizzare la topologia delle decisioni.
L'applicazione di metriche geometriche, quali la Cosine Similarity e la Distanza Euclidea, ha permesso la corretta interpretazione degli attacchi testuali nello spazio latente: contrariamente a quanto si potesse ipotizzare, gli attacchi testuali non spingono il vettore lungo l'asse della rappresentazione dei concetti di sicurezza e vulnerabilità della baseline, bensì provocano uno shift violento della traiettoria, deviando lo stato latente in dimensioni ortogonali. Il modello viene ingannato non perché riconosce il codice come sicuro, ma perché gli attacchi Black-Box distruggono la visione del task originario costringendo la rete ad operare in una zona grigia. L'attacco White-Box di Activation Steering si è dimostrato fedele alle ipotesi iniziali solo al momento dell'iniezione: nel layer in cui avviene l'attacco lo spostamento viene effettuato sull'asse del vettore di steering per poi allontanarsi nei livelli successivi.

L'indagine condotta con l'attacco di Logit Bias e la tecnica di Contrastive Logit Lens ha permesso di delineare una discrepanza tra la divergenza geometrica e la decodifica finale: se le misurazioni spaziali tramite la norma L2 mostrano un progressivo allontanamento dal vettore di riferimento, l'analisi dei logit rivela che nei layer finali avviene un tentativo di autocorrezione, cercando di preservare la bussola logica della vulnerabilità. Gran parte degli attacchi ha compiuto un dirottamento completo con l'impossibilità di recupero, mentre una piccola percentuale di successi evidenzia come le perturbazioni introdotte abbiano assottigliato i margini di confidenza e sconvolto le probabilità associate a varianti ortografiche o sintattiche, portando al collasso dell'output a schermo.
Ciò evidenzia non solo la naturale evoluzione della computazione del modello, ma conferma le ipotesi di ripartizione dei compiti tra i layer: la vera decodifica logica della vulnerabilità avviene, e viene compromessa, nei layer centrali. Ai livelli finali viene lasciata l'elaborazione sintattica della risposta dove, operando su un vettore ormai corrotto, il modello non possiede la forza strutturale per invertire il dirottamento.

L'adozione di un modello con capacità di ragionamento differito assieme alla modifica del prompt ha portato alla luce le vere potenzialità della Chain-of-Thought: il \textit{Thinking} del modello, sia indotto che nativo, ha portato ad una resistenza maggiore agli attacchi.

\section{Sviluppi Futuri}
Lo studio si propone di essere un punto di partenza per ricerche successive: sebbene in letteratura si sostenga che il numero di parametri di un LLM non influenzi la corretta identificazione di perturbazioni, risulta metodologicamente corretto avviare una ricerca adottando modelli con un numero di parametri più elevato rispetto a quelli presentati, oltre ad estendere la revisione a codici più lunghi per poter analizzare come l'attacco si sviluppa.
Inoltre, l'utilizzo di soli tre attacchi Black-Box non permette la completa visione dei parallelismi delle perturbazioni introdotte. Ampliare lo studio della cosine similarity a più attacchi di natura Black-Box aiuterebbe a mappare lo spettro su cui si collocano, permettendo in futuro di sviluppare tecniche di difesa o collaudo nella sicurezza dei modelli stessi.
Dato l'ampio uso dei vettori di steering e la mappatura dello spazio latente, è necessario provare ad implementare meccaniche di positive steering, già presenti in letteratura, affiancate a sentinelle nello spazio latente per una migliore correzione del pensiero del modello.